% =====================================================================
%  Beneath the Mesa -- ENG 220 Project Report (Team 23)
%  Two-column report template styled to match the project proposal:
%  Palatino (URW Palladio) text, navy headings, blue rules and header.
%
%  Build:  pdflatex report -> bibtex report -> pdflatex report (x2)
%          (or `latexmk -pdf report.tex`; Overleaf does this automatically)
%
%  Orange [notes] mark what still needs to be written.
%  Dashed boxes mark figures that still need to be made.
% =====================================================================
\documentclass[10pt,twocolumn,letterpaper]{article}

% ---------------------------------------------------------------------
%  Switches
% ---------------------------------------------------------------------
\newif\ifdraftnotes \draftnotestrue   % \draftnotesfalse hides every [orange note]
\newif\ifcoverpage  \coverpagefalse   % \coverpagetrue adds a proposal-style cover page
\newif\ifroughdraft \roughdrafttrue    % \roughdraftfalse removes the red ROUGH DRAFT stamp

% ---------------------------------------------------------------------
%  Fonts and text
% ---------------------------------------------------------------------
\usepackage[T1]{fontenc}
\usepackage[utf8]{inputenc}
\usepackage{mathpazo}                 % URW Palladio, the proposal's typeface
\usepackage[protrusion=true,expansion=true]{microtype}
\usepackage{amsmath}

% ---------------------------------------------------------------------
%  Page layout
% ---------------------------------------------------------------------
\usepackage[letterpaper,top=0.8in,bottom=0.85in,left=0.75in,right=0.75in,
            columnsep=0.28in,headheight=20pt,headsep=8pt,footskip=26pt]{geometry}
\setlength{\parindent}{1.2em}
\setlength{\parskip}{0pt}

% ---------------------------------------------------------------------
%  Colors (sampled from the proposal PDF)
% ---------------------------------------------------------------------
\usepackage{xcolor}
\definecolor{titleblue}{RGB}{23,54,93}     % #17365D  headings, title
\definecolor{accentblue}{RGB}{54,95,145}   % #365F91  rules, header, links
\definecolor{tintblue}{RGB}{237,241,247}   % light panel fill
\definecolor{draftorange}{RGB}{190,84,20}  % draft notes only
\definecolor{draftred}{RGB}{192,0,0}       % ROUGH DRAFT stamp

% ---------------------------------------------------------------------
%  Headings
% ---------------------------------------------------------------------
\usepackage{titlesec}
\usepackage{needspace}
\usepackage{etoolbox}
\titleformat{\section}{\color{titleblue}\large\bfseries}{\thesection}{0.55em}{}
  [{\color{accentblue}\titlerule[0.6pt]}]
\titlespacing*{\section}{0pt}{2.0ex plus 0.6ex minus 0.3ex}{1.0ex}
\titleformat{\subsection}{\color{accentblue}\normalsize\bfseries}{\thesubsection}{0.55em}{}
\titlespacing*{\subsection}{0pt}{1.4ex plus 0.4ex minus 0.2ex}{0.5ex}
\titleformat{\subsubsection}[runin]{\normalsize\bfseries\itshape}{}{0pt}{}[.]
\titlespacing*{\subsubsection}{0pt}{1.0ex plus 0.3ex}{0.6em}
\pretocmd{\section}{\Needspace{4\baselineskip}}{}{}       % keep headings with their text
\pretocmd{\subsection}{\Needspace{3\baselineskip}}{}{}

% ---------------------------------------------------------------------
%  Running header and footer (same as the proposal)
% ---------------------------------------------------------------------
\usepackage{fancyhdr}
\pagestyle{fancy}
\fancyhf{}
\fancyhead[L]{\small\color{accentblue}Project Report}
\fancyhead[R]{\small\color{accentblue}Aquifer, Well \& Acequia Resources: Sandoval County, NM}
\fancyfoot[C]{\small\thepage}
\renewcommand{\headrulewidth}{0.4pt}
\renewcommand{\headrule}{\hbox to\headwidth{\color{accentblue}\leaders\hrule height\headrulewidth\hfill}}

\fancypagestyle{firstpage}{%
  \fancyhf{}%
  \renewcommand{\headrulewidth}{0pt}%
  \fancyfoot[L]{\footnotesize\color{accentblue}ENG 220 Project Report \,\textbullet\, Team 23 \,\textbullet\, \reportstatus}%
  \fancyfoot[R]{\small\thepage}%
}

% ---------------------------------------------------------------------
%  Lists, tables, figures, captions
% ---------------------------------------------------------------------
\usepackage{enumitem}
\setlist{itemsep=1pt,topsep=3pt,parsep=0pt}
\setlist[itemize]{leftmargin=1.2em,label={\color{accentblue}\textbullet}}
\usepackage{booktabs,tabularx,array}
\renewcommand{\arraystretch}{1.12}
\usepackage{graphicx}
\graphicspath{{figures/}}
\usepackage{caption}
\captionsetup{font=small,labelfont={bf,color=titleblue},labelsep=period,justification=justified}
\captionsetup[table]{position=top,skip=4pt}  % for tables you add later
\usepackage{tikz}
\usepackage{stfloats}                 % lets full-width figures sit at the page bottom too

% ---------------------------------------------------------------------
%  References and links
% ---------------------------------------------------------------------
\usepackage[round,authoryear]{natbib}
\setlength{\bibsep}{2.5pt}
\renewcommand{\bibfont}{\small\raggedright}
\usepackage{xurl}
\usepackage[colorlinks=true,linkcolor=accentblue,citecolor=accentblue,urlcolor=accentblue,
            pdftitle={Beneath the Mesa: Aquifer, Well, and Acequia Water Resources Across Sandoval County, New Mexico},
            pdfauthor={Djordjevic, Gonzales, Wanueka, Watkins, Yousef}]{hyperref}
\urlstyle{same}

% ---------------------------------------------------------------------
%  Draft helpers
% ---------------------------------------------------------------------
\newcommand{\reportstatus}{\ifroughdraft Rough draft, October 25, 2026\else Final, November 30, 2026\fi}
% Red stamp shown on page 1 and the cover page while \roughdrafttrue
\newcommand{\draftstamp}[1][\large]{\ifroughdraft%
  \tikz\node[draw=draftred,line width=1.1pt,rounded corners=2pt,inner xsep=9pt,inner ysep=4pt,
    text=draftred,font=#1\bfseries]{\textls[180]{ROUGH DRAFT}};\fi}
\newcommand{\draftnote}[1]{\ifdraftnotes{\color{draftorange}\itshape[#1]}\fi}
% Dashed box that stands in for a figure that does not exist yet.
%   \figplaceholder[height]{what the figure will show}
% When the figure is ready, save it in figures/ and replace the
% placeholder with:  \includegraphics[width=\linewidth]{file-name}
\newcommand{\figplaceholder}[2][1.8in]{%
  \begin{tikzpicture}
    \draw[accentblue!70,dashed,line width=0.6pt,fill=tintblue,rounded corners=2pt]
      (0,0) rectangle (\linewidth,#1);
    \node[text width=0.82\linewidth,align=center,font=\small\itshape,text=accentblue]
      at (0.5\linewidth,0.5*#1) {#2};
  \end{tikzpicture}}
\newcommand{\infohead}[1]{{\color{titleblue}\footnotesize\bfseries\MakeUppercase{#1}}\par\vspace{2pt}}

% =====================================================================
\begin{document}

% ---------------------------------------------------------------------
%  Optional cover page (same layout as the proposal cover)
% ---------------------------------------------------------------------
\ifcoverpage
\begin{titlepage}
  \centering
  \vspace*{1.6in}
  {\color{accentblue}\rule{\textwidth}{0.8pt}}\par\vspace{1.4em}
  {\color{titleblue}\fontsize{24}{29}\selectfont\bfseries Beneath the Mesa\par}
  \vspace{0.9em}
  {\color{accentblue}\fontsize{17}{22}\selectfont Assessing Aquifer, Well, and Acequia Water\\
    Resources Across Sandoval County, New Mexico\par}
  \vspace{1.3em}
  {\color{accentblue}\rule{\textwidth}{0.8pt}}\par\vspace{1.6em}
  {\itshape A County-Scale Screening Study with Cuba, New Mexico as a Focus Area\par}
  \vspace{1.3in}
  {\color{titleblue}\large\bfseries Project Report\par}
  \ifroughdraft\vspace{1.2em}\draftstamp[\Large]\par\vspace{-0.4in}\fi
  \vspace{1.2in}
  \begin{tabular}{@{}r@{\hspace{0.8em}}l@{}}
    \textbf{Course:}     & ENG 220 \\
    \textbf{Instructor:} & Ramiro Jordan \\
    \textbf{Team:}       & \#23 Petar Djordjevic, Dylan Gonzales, James Wanueka,\\
                         & \phantom{\#23 }Micah Watkins, Meckrem Yousef \\
    \textbf{Date:}       & \ifroughdraft October 25, 2026 (rough draft)\else November 30, 2026\fi \\
  \end{tabular}
  \vfill
\end{titlepage}
\fi

% ---------------------------------------------------------------------
%  Title block and abstract (full width)
% ---------------------------------------------------------------------
\twocolumn[{%
\begin{center}
  \ifroughdraft\vspace*{-1.5ex}\draftstamp\par\vspace{1.4ex}\fi
  {\color{accentblue}\rule{\textwidth}{0.8pt}}\par\vspace{0.9ex}
  {\color{titleblue}\fontsize{22}{26}\selectfont\bfseries Beneath the Mesa\par}
  \vspace{0.7ex}
  {\color{accentblue}\fontsize{14.5}{18}\selectfont Assessing Aquifer, Well, and Acequia Water Resources\\
    Across Sandoval County, New Mexico\par}
  \vspace{0.9ex}
  {\color{accentblue}\rule{\textwidth}{0.8pt}}\par\vspace{0.9ex}
  {\itshape A County-Scale Screening Study with Cuba, New Mexico as a Focus Area\par}
  \vspace{2.0ex}
  {\large Petar Djordjevic, Dylan Gonzales, James Wanueka, Micah Watkins, and Meckrem Yousef\par}
  \vspace{0.6ex}
  {\small ENG 220, Team 23 \,\textbullet\, Instructor: Ramiro Jordan\par}
\end{center}
\vspace{1.2ex}
{\color{accentblue}\rule{\textwidth}{0.4pt}}\par\vspace{1.4ex}
\noindent
\begin{minipage}[t]{0.25\textwidth}
  \raggedright\footnotesize
  \infohead{Report info}
  \ifroughdraft Status: {\color{draftred}\bfseries Rough draft}\\
  Date: October 25, 2026\\
  \else Status: Final\\
  \fi Final due: November 30, 2026\par\vspace{1.4ex}
  \infohead{Keywords}
  Groundwater\\ Aquifers\\ Acequias\\ Sandoval County\\ Cuba, New Mexico\par\vspace{1.4ex}
  \infohead{Online}
  \href{https://eng220aquiferproject.streamlit.app/}{Interactive dashboard}\\
  \href{https://github.com/pdjordj11/eng_220_aquifer_project}{Code and data (GitHub)}\\
  \href{https://pdjordj11.github.io/eng_220_aquifer_project/}{Project documentation}
\end{minipage}%
\hfill{\color{accentblue}\vrule width 0.4pt}\hfill
\begin{minipage}[t]{0.70\textwidth}
  \infohead{Abstract}
  \small
  Water scarcity is a persistent challenge across New Mexico. Sandoval County spans several
  hydrogeologic settings, and its communities depend on a mix of groundwater wells, surface water,
  and acequia irrigation systems. This project evaluates the condition and long-term viability of
  those water resources across the county, with Cuba, New Mexico, on the eastern edge of the San
  Juan Basin, as a focused case study. The project is motivated by concern about declining
  groundwater near Cuba, but rather than assume how severe that decline is or what is causing it,
  we test what the available public data actually show.
  \draftnote{Add the data used, the main findings, and what they mean, in three to four
  sentences, once the analysis is complete.}
\end{minipage}
\par\vspace{1.6ex}
{\color{accentblue}\rule{\textwidth}{0.4pt}}\par\vspace{2.6ex}
}]
\thispagestyle{firstpage}

% =====================================================================
\section{Introduction}
% =====================================================================
\subsection{Background and Motivation}
\draftnote{Background on water resources in New Mexico and Sandoval County, and the concern about
declining groundwater near Cuba that motivated the project.}

\subsection{Approach}
The study works at two levels:
\begin{itemize}
  \item \textbf{County-wide screening.} Map and compare aquifers, wells, water levels, water
        quality, and acequia and surface-water features across Sandoval County. Because the county
        crosses several geologic regions, comparisons are grouped by region or aquifer system
        instead of treating the county as one uniform aquifer.
  \item \textbf{Cuba case study.} A closer look at the Ojo Alamo Sandstone and San Jose Formation
        near Cuba, using San Juan Basin aquifer datasets alongside nearby well, water-level, and
        water-quality records.
\end{itemize}

\subsection{Research Questions}
\begin{enumerate}[label={\color{titleblue}\bfseries RQ\arabic*.},leftmargin=2.6em,labelsep=0.4em]
  \item How do groundwater levels vary across Sandoval County, and which areas show sustained
        decline, stability, or recovery?
  \item What do well records show about well depth, depth to water, and the aquifer units that
        communities rely on?
  \item Do the available data support the concern that Cuba's main groundwater source is
        approaching severe depletion?
  \item How do the Ojo Alamo Sandstone and San Jose Formation compare near Cuba in extent,
        thickness, hydraulic properties, water levels, and water quality?
  \item How does groundwater quality, including salinity, vary across the county and near Cuba?
  \item Where are acequias located, what surface water do they depend on, and which communities
        do they serve?
  \item Do changes in groundwater levels line up with nearby surface-water or acequia conditions,
        or are the records too limited to tell?
  \item Where is monitoring coverage strong, and where are the most important data gaps?
  \item Which indicators can be combined into a dashboard that shows water-resource stress without
        overstating how certain the data are?
\end{enumerate}

% ---------------------------------------------------------------------
%  Figure 1: county map (LaTeX places it at the top of the next page)
% ---------------------------------------------------------------------
\begin{figure*}[t]
  \centering
  \figplaceholder[2.2in]{Figure 1 (to be added by the analysis lead): county map showing wells,
  water-level measurements, major aquifer and geologic regions, and acequias, with Cuba marked.}
  \caption{\draftnote{Caption to be written.}}
  \label{fig:map}
\end{figure*}

% =====================================================================
\section{Study Area}\label{sec:study}
% =====================================================================
\subsection{Sandoval County}
\draftnote{Describe the county, its hydrogeologic settings, and its main communities and water
sources (Fig.~\ref{fig:map}).}

\subsection{Cuba Focus Area}
Cuba lies on the eastern edge of the San Juan Basin. The case study focuses on the Ojo Alamo
Sandstone and San Jose Formation near Cuba \citep{usgs_rasa_sjb,kernodle1996}.
\draftnote{Describe Cuba, its water supply, and the aquifers it draws from, and define the
boundaries of the focus area.}

% =====================================================================
\section{Data}\label{sec:data}
% =====================================================================
Town-specific data for Cuba are limited, so the study relies on broader public datasets that can
be filtered geographically.

\subsection{Datasets Used}
\subsubsection{USGS monitoring statistics}
USGS monthly and yearly statistics for 30 monitoring sites in Sandoval County (17 stream gages,
9 wells, 2 canals, and 2 reservoirs), with records from 1926 to 2026. The file covers streamflow,
depth to groundwater, reservoir storage, suspended sediment, water temperature, and water-quality
measures: specific conductance, dissolved oxygen, pH, and turbidity.
\draftnote{Source to be added.}

\subsubsection{OSE aquifer tests}
New Mexico Office of the State Engineer aquifer-test records, 3,455 statewide and 179 in Sandoval
County, including well depth, screened interval, geologic formation, and hydraulic properties
such as transmissivity and storativity.
\draftnote{Source to be added.}

\subsection{Datasets to Be Collected}
\begin{itemize}
  \item \textbf{Aquifer geometry near Cuba:} USGS San Juan Basin digital aquifer data for the Ojo
        Alamo Sandstone and San Jose Formation \citep{usgs_rasa_sjb}.
  \item \textbf{Water quality and salinity:} the San Juan Basin Brackish Water Assessment
        \citep{nmbg_bwa} and New Mexico Water Data records \citep{nmwaterdata}.
  \item \textbf{Acequias:} locations and service areas of acequias and community ditches
        \citep{isc_acequias}.
  \item \textbf{More water-level records:} the National Ground-Water Monitoring Network
        \citep{ngwmn} and the Healy Collaborative Groundwater Monitoring Network \citep{cgmn}.
  \item \textbf{County boundary:} a Sandoval County boundary layer to organize the maps.
\end{itemize}
\draftnote{Data acquisition lead: explain why each dataset was chosen and which research questions
it answers.}

% =====================================================================
\section{Methods}\label{sec:methods}
% =====================================================================
\subsection{Data Cleaning}
Raw downloads are kept unedited in the repository. Each cleaned file is made by a script that
reads the raw file and writes the cleaned one, so anyone can rerun it and get the same result.
All cleaned files share the same columns for site ID, site name, latitude, longitude, date, and
source, with one row per measurement and missing values left blank.
\draftnote{Data cleaning lead: describe each cleaning step, the decisions made, and the final row
counts.}

\subsection{Analysis}
The analysis focuses on measurable indicators, such as depth to water, water-level change over
time, well depth, aquifer unit, and salinity, rather than a full hydrologic model. Where records
are sparse, the results say so instead of filling the gaps.
\draftnote{Analysis lead: describe the method used for each research question.}

\subsection{Classifying Findings}
Each finding is labeled as one of the following:
\begin{itemize}
  \item \textbf{Measured:} a trend or pattern the data show directly.
  \item \textbf{Hypothesis:} a possible explanation the data are consistent with but don't prove.
  \item \textbf{Insufficient data:} a question the available records can't answer yet.
\end{itemize}

\subsection{Interactive Dashboard}
Results are presented in an interactive Streamlit dashboard built with pandas and Plotly. Code
and documentation are maintained on GitHub \citep{team23_repo}.
\draftnote{Dashboard lead: describe each page of the app and what a visitor can do there.}

% =====================================================================
\section{Results}\label{sec:results}
% =====================================================================
\subsection{Groundwater Levels}
\begin{figure}[t]
  \centering
  \figplaceholder[1.25in]{Figure 2 (to be added): time series of groundwater levels for wells with
  repeated measurements.}
  \caption{\draftnote{Caption to be written.}}
  \label{fig:levels}
\end{figure}
\draftnote{RQ1 (Fig.~\ref{fig:levels}).}

\subsection{Wells and Aquifer Units}
\draftnote{RQ2.}

\subsection{Cuba Case Study}
\draftnote{RQ3 and RQ4.}

\subsection{Water Quality}
\draftnote{RQ5.}

\subsection{Acequias and Surface Water}
\draftnote{RQ6 and RQ7.}

\subsection{Monitoring Coverage and Data Gaps}
\draftnote{RQ8.}

% =====================================================================
\section{Discussion}\label{sec:discussion}
% =====================================================================
\draftnote{What the results mean for Sandoval County and Cuba, and which indicators best
communicate water-resource stress (RQ9).}

% =====================================================================
\section{Limitations}\label{sec:limits}
% =====================================================================
\draftnote{Data gaps and what the data can't support.}

% =====================================================================
\section{Conclusions}\label{sec:conclusions}
% =====================================================================
\draftnote{Summary of findings and recommended next steps.}

% =====================================================================
%  Back matter
% =====================================================================
\Needspace{9\baselineskip}
\section*{Data and Code Availability}
\begin{description}[style=nextline,leftmargin=1em,itemsep=1pt,font=\normalfont\bfseries]
  \item[Code and data] \href{https://github.com/pdjordj11/eng_220_aquifer_project}{\nolinkurl{github.com/pdjordj11/eng_220_aquifer_project}}
  \item[Interactive dashboard] \href{https://eng220aquiferproject.streamlit.app/}{\nolinkurl{eng220aquiferproject.streamlit.app}}
  \item[Documentation] \href{https://pdjordj11.github.io/eng_220_aquifer_project/}{\nolinkurl{pdjordj11.github.io/eng_220_aquifer_project}}
\end{description}

\section*{Author Contributions}
\begin{description}[leftmargin=0pt,labelsep=0.4em,itemsep=0pt,font=\normalfont\bfseries]
  \item[P.~Djordjevic:] documentation, final report, presentation, and video.
  \item[D.~Gonzales:] data acquisition and dataset justification.
  \item[J.~Wanueka:] data cleaning and preprocessing.
  \item[M.~Watkins:] analysis code and visualizations.
  \item[M.~Yousef:] GitHub workflow and the Streamlit app.
\end{description}

\section*{Acknowledgments}
\draftnote{To be written.}

\bibliographystyle{plainnat}
\bibliography{references}

\end{document}
